\documentclass[11pt]{article}
\usepackage[margin=1in]{geometry}
\usepackage{amsmath,amssymb,amsthm}
\usepackage{array}
\usepackage{parskip}
\usepackage{booktabs}

\newtheorem{theorem}{Theorem}
\newtheorem{lemma}[theorem]{Lemma}
\newtheorem{corollary}[theorem]{Corollary}
\newtheorem{proposition}[theorem]{Proposition}
\theoremstyle{definition}
\newtheorem{definition}[theorem]{Definition}
\theoremstyle{remark}
\newtheorem{remark}[theorem]{Remark}

\newcommand{\TM}{\mathrm{TM}}
\newcommand{\code}[1]{\texttt{#1}}

\emergencystretch=3em

\title{Disproof of conjecture \texttt{00000001260}}
\author{earthking11}
\date{2026-09-13}

\begin{document}
\maketitle

\section*{Verdict: FALSE}

We disprove conjecture \texttt{00000001260} as stated. The conjecture asserts a
\emph{universal} upper bound, $p(n) \le n+2$, for the factor complexity of
every primitive substitution in the ``purely substitutive class''. The
Thue--Morse word, the fixed point of the primitive substitution
$0 \mapsto 01$, $1 \mapsto 10$, already violates this bound at $n = 3$, where
$p(3) = 6 > 5 = 3+2$. The failure is not an isolated accident: the celebrated
Tribonacci word named in the conjecture itself satisfies $p(n) = 2n+1 > n+2$
for every $n \ge 2$, so it too refutes the stated bound.

\section{The conjecture}

Quoted from \code{conjectures/00000001260.md}:

\begin{quote}
\textbf{Definition:} The factor complexity $p(n)$ of a primitive substitution
counts distinct factors of length $n$. \textbf{Conjecture:} The optimal upper
bound is $p(n) \le n+2$ within the purely substitutive class, attained by
three-letter substitution families beyond the Tribonacci word; $p(n) = n+1$
occurs only for Sturmian words.
\end{quote}

The filed statement consists of a \emph{definition} (factor complexity), a
\emph{universal bound} ($p(n) \le n+2$ over the class), an \emph{attainment}
claim (three-letter families attain it), and a \emph{classification} claim
($p(n)=n+1$ only for Sturmian words). The universal bound is the load-bearing
part; we refute it, and thereby the conjecture. The second clause,
$p(n)=n+1$ for Sturmian words, is true and we confirm it below.

\section{Definitions}

\begin{definition}[Substitution and fixed point]
A \emph{substitution} is a morphism $\sigma$ on words over a finite alphabet.
A word $w$ is a \emph{fixed point} of $\sigma$ if $\sigma(w)=w$ and $w$ begins
with $\sigma(w_0)$; the fixed point is obtained as
$\lim_{k\to\infty}\sigma^k(0)$ when $\sigma(0)$ begins with $0$.
\end{definition}

\begin{definition}[Factor complexity]
For an infinite word $w$, the \emph{factor complexity} $p(n)$ is the number of
distinct finite words of length $n$ (its \emph{factors}) occurring as a
contiguous block in $w$.
\end{definition}

\begin{definition}[Primitive substitution]
A substitution $\sigma$ on a $d$-letter alphabet is \emph{primitive} if there
is $k \ge 1$ such that every entry of the incidence matrix $M^k$ is strictly
positive, where $M_{ij}$ is the number of occurrences of letter $j$ in
$\sigma(i)$.
\end{definition}

\begin{definition}[Thue--Morse word]
The Thue--Morse substitution on $\{0,1\}$ is
\[
\sigma_{\TM}\colon\quad 0 \mapsto 01,\qquad 1 \mapsto 10 .
\]
Its fixed point starting with $0$ is the \emph{Thue--Morse word}
$\TM = 0110100110010110\cdots$. Equivalently, the $n$-th letter of $\TM$
is the parity of the binary digit sum of $n$.
\end{definition}

\section{The Thue--Morse counterexample}

\begin{proposition}[Primitivity]
$\sigma_{\TM}$ is primitive. Its incidence matrix is
$M = \begin{pmatrix} 1 & 1 \\ 1 & 1 \end{pmatrix}$, every entry of which is
positive; hence $M^k$ is positive for every $k \ge 1$, so the defining
condition holds already with $k=1$.
\end{proposition}

\begin{proof}
Immediate from the definition: $\sigma_{\TM}(0) = 01$ contains one $0$ and one
$1$, and likewise $\sigma_{\TM}(1) = 10$. Thus $M$ has all entries $1>0$.
\end{proof}

The Thue--Morse word is standardly \emph{purely substitutive}: it is the fixed
point of an actual substitution, and this substitution is primitive. It is a
constant-length (of length $2$) substitution, and the filed conjecture text
imposes no restriction excluding constant-length substitutions and no
restriction on the size of the alphabet.

\begin{theorem}[The bound fails at $n=3$]\label{thm:main}
The Thue--Morse word has at least six distinct factors of length $3$; indeed
\[
p_{\TM}(3) = 6 > 5 = 3 + 2 .
\]
Consequently the asserted universal bound $p(n) \le n+2$ is false within the
purely substitutive class.
\end{theorem}

\begin{proof}
The six blocks
\[
011,\quad 110,\quad 101,\quad 010,\quad 100,\quad 001
\]
each occur in $\TM$ as a length-$3$ factor: they start at positions
$0,1,2,3,4,5$ of the prefix $0110100110010110\cdots$, respectively. They are
pairwise distinct, so $p_{\TM}(3) \ge 6$. A direct enumeration over the prefix
of length $2^{18}$ gives exactly $6$ distinct length-$3$ factors, so
$p_{\TM}(3)=6$. Since $6 > 5$, the bound $p(n) \le n+2$ fails at $n=3$.
\end{proof}

\begin{corollary}
Conjecture \texttt{00000001260} is false: its universal upper bound is
violated by a primitive, purely substitutive word.
\end{corollary}

\begin{remark}[Only a lower bound is needed]
To refute a universal \emph{upper} bound it suffices to exhibit, for a single
word, a single $n$ and a \emph{lower} bound $p(n) > n+2$. Each of the six
blocks above is a genuine factor by construction (it starts at an explicit
position), so the count $6$ is a lower bound on $p_{\TM}(3)$; no completeness
argument (no proof that there are \emph{at most} six) is required. This is
exactly what the Lean formalisation records.
\end{remark}

\section{The complexity table for $n = 1,\dots,15$}

Counting distinct factors of $\TM$ over the prefix of length $2^{18}$:

\begin{center}
\begin{tabular}{rrrrrrrrrrrrrrrr}
\toprule
$n$   & 1 & 2 & 3 & 4  & 5  & 6  & 7  & 8  & 9  & 10 & 11 & 12 & 13 & 14 & 15\\
\midrule
$p(n)$ & 2 & 4 & 6 & 10 & 12 & 16 & 20 & 22 & 24 & 28 & 32 & 36 & 40 & 42 & 44\\
$n+2$  & 3 & 4 & 5 & 6  & 7  & 8  & 9  & 10 & 11 & 12 & 13 & 14 & 15 & 16 & 17\\
\bottomrule
\end{tabular}
\end{center}

The bound $p(n) \le n+2$ holds only at $n = 2$ (where $p(2) = 4 = 2+2$) and
fails for every $n \ge 3$, the first failure being $p(3) = 6 > 5$. Numerical
verification confirms $p(n) > n+2$ for every $n \in [3,300]$.

\section{Bonus counterexamples}

\subsection{The Tribonacci word (the word the conjecture names)}

The Tribonacci word is the fixed point of the primitive substitution
\[
0 \mapsto 01,\qquad 1 \mapsto 02,\qquad 2 \mapsto 0 .
\]
It is a three-letter substitution, and it is precisely the ``Tribonacci word''
the conjecture refers to. Its factor complexity is known to be exactly linear,
\[
p(n) = 2n+1 \qquad \text{for all } n \ge 1 .
\]
Hence $p(n) = 2n+1 > n+2$ for every $n \ge 2$. So the very word named in the
conjecture violates the asserted universal bound, not merely some exotic
example. Computation over a prefix of length $223317$ reproduces
$p(1..15) = 3,5,7,9,11,13,15,17,19,21,23,25,27,29,31$, matching $2n+1$.

\subsection{The period-doubling word}

The period-doubling word is the fixed point of
\[
0 \mapsto 01,\qquad 1 \mapsto 00 .
\]
Computing over a prefix of length $65536$ gives
\[
p(1),p(2),p(3),p(4),p(5),\dots = 2,3,5,6,8,\dots
\]
so the bound $p(n) \le n+2$ first fails at $n = 5$, where $p(5) = 8 > 7$.

\subsection{Sturmian words confirm the second clause}

The Fibonacci word, fixed point of $0 \mapsto 01$, $1 \mapsto 0$, is the
archetypal Sturmian word, and it satisfies $p(n) = n+1$ for all $n \ge 1$.
Computation over a prefix of length $131072$ gives
$p(1..14) = 2,3,4,\dots,15$, exactly $n+1$. Thus the conjecture's
classification clause ($p(n)=n+1$ only for Sturmian words) is consistent with
our data; it is the universal bound that fails.

\section{Reading caveat (honesty statement)}

The refutation above targets the \emph{literal universal reading} of the filed
text: ``the optimal upper bound is $p(n) \le n+2$ within the purely
substitutive class'', with no alphabet restriction and no exclusion of
constant-length substitutions.

\begin{itemize}
\item If ``purely substitutive'' were silently restricted to \emph{three-letter}
substitutions, or to \emph{Pisot-type} substitutions, the Thue--Morse
counterexample (a two-letter, constant-length substitution) would be excluded
and the refutation would not apply. \emph{No such restriction appears in the
filed text.}
\item If the conjecture were reinterpreted as a statement about the
\emph{minimal growth rate over the class} rather than a universal upper bound,
then the Thue--Morse word would again be irrelevant. Again, \emph{no such
reinterpretation is present in the text}: the sentence plainly asserts a bound
$p(n) \le n+2$ holding within the class.
\item The sentence itself is garbled: it mixes a universal bound
($p(n) \le n+2$ ``within the class'') with an \emph{attainment} claim
(``attained by three-letter substitution families beyond the Tribonacci
word''), which are statements of different logical type. It is not a precise
mathematical claim, and no repaired reading is supplied by the text.
\end{itemize}

\noindent Under the literal reading there is no ambiguity: the bound is false.

\section{Reproducibility}

The checks are carried out by \code{reproduce.py}, which uses the Python~3
standard library only. It builds the Thue--Morse prefix of length $2^{18}$ by
iterated substitution, counts distinct length-$n$ factors for $n=1,\dots,300$,
verifies $p(3)=6$ and $p(n) > n+2$ for all $n \in [3,300]$, and likewise
computes the period-doubling, Tribonacci, and Fibonacci complexities:

\begin{quote}\ttfamily
\$ python3 reproduce.py\\
... \textit{(tables and checks)}\\
PASS: conjecture 00000001260 is FALSE.
\end{quote}

It prints \code{PASS} and exits $0$ exactly when every check holds.

\section{Formalisation}

The core of the disproof is formalised in the accompanying Lean~4 project
under \code{lean4/}, built with \code{lake build} and audited by
\code{lake env lean Check.lean}. It uses core Lean only (\code{import Std}, no
Mathlib) and no \code{sorry}. The Thue--Morse word is defined by a structural
recursion so that the concrete computation reduces under \code{decide}, and
the key statements are:

\begin{quote}\ttfamily
theorem tm\_p3\_ge\_6 : 6 \(\le\) fac3.eraseDups.length\\
theorem tm\_p3\_eq\_6 : fac3.eraseDups.length = 6\\
theorem tm\_refutes\_bound : fac3.eraseDups.length \(>\) 3 + 2\\
theorem tm\_prefix\_agrees : (List.range 64).map tm = gen 6\\
theorem incidence\_positive : \(\forall\) r \(\in\) incidence, \(\forall\) c \(\in\) r, 0 \(<\) c\\
theorem conjecture\_00000001260\_false :\\
\hspace*{2em}fac3.eraseDups.length = 6 \(\wedge\) \(\neg\) (fac3.eraseDups.length \(\le\) 3 + 2)
\end{quote}

\noindent The formalisation records exactly the lower bound $p(3) \ge 6$:
\code{fac3} is the list of length-$3$ blocks starting at explicit positions
$0,\dots,61$, so each of its entries is a genuine factor, and
\code{fac3.eraseDups.length} is therefore a lower bound for $p(3)$. The
consistency check \code{tm\_prefix\_agrees} ties the closed-form definition of
$\TM$ to the substitution-generated prefix, and \code{incidence\_positive}
witnesses primitivity of $0 \mapsto 01$, $1 \mapsto 10$.

\begin{sloppypar}
\code{Check.lean} prints \code{\#print axioms} for every theorem and reports
\code{propext} only; \code{incidence\_positive} depends on no axioms at all.
In particular it reports neither \code{sorryAx} nor
\code{Lean.ofReduceBool}; no Mathlib is involved.
\end{sloppypar}

\section*{Files}

\begin{tabular}{@{}p{0.26\textwidth}p{0.70\textwidth}@{}}
\toprule
File & Description\\
\midrule
\code{README.md} & Verdict, bilingual quote, definitions, tables, caveat, status.\\
\code{main.tex} & This source; standalone \code{article}, \code{tectonic --outdir build main.tex}.\\
\code{build/main.pdf} & Compiled PDF.\\
\code{reproduce.py} & Stdlib-only reproduction; prints \code{PASS}/\code{FAIL}.\\
\code{lean4/} & Lean 4 project: \code{Main.lean}, \code{Check.lean}, \code{lakefile.toml}, \code{lean-toolchain}, \code{README.md}.\\
\bottomrule
\end{tabular}

\end{document}
